\documentclass[a4paper,11pt,titlepage,singlespacing]{article}
\usepackage[T1]{fontenc} 
\usepackage[utf8]{inputenc}
\usepackage[brazilian]{babel}

\usepackage{amsmath}
\usepackage{indentfirst}
\usepackage{graphicx}
\usepackage{hyperref}
\usepackage{multicol,lipsum}
\usepackage[bottom=2cm,top=2cm,left=2cm,right=2cm]{geometry}
%\usepackage{color}
\usepackage{cancel}
\usepackage{soulutf8}
\usepackage{color}

%%%% CALUÃ COMEÇOU A MEXER DAQUI %%%%%

\usepackage[easyscsl]{kpfonts}
\usepackage[T1]{fontenc}

\usepackage{floatrow}
{\renewcommand{\arraystretch}{1.25}
\usepackage{arydshln}

\newcommand\todo[1]{\textcolor{red}{#1}}

% \usepackage[hidelinks]{hyperref}
\usepackage[alf,versalete]{abntex2cite}

\renewcommand{\thesection}{\oldstylenums{\arabic{section}}}
\renewcommand{\thesubsection}{\oldstylenums{\arabic{section} · \arabic{subsection}}}
\renewcommand{\thepage}{\oldstylenums{\arabic{page}}}

\renewcommand{\yearstyle}{\oldstylenums}

\linespread{1.25}



\setlength{\topmargin}{-0.5cm} % extra vert. space + at the top of header: 23pt -15
\setlength{\oddsidemargin}{0.5cm} % extra spc added at the left of odd page: 0pt
% \setlength{\evensidemargin}{1cm} % ext. spc added at the left of even pg: 59pt
\setlength{\textheight}{23.5cm} % height of the body: 592pt
\setlength{\textwidth}{15cm}
\newcommand{\ignore}[1]{}


% https://tex.stackexchange.com/questions/255167/hanging-indent-in-headings
\usepackage{titlesec}
\newlength\TitleOverhang
\setlength\TitleOverhang{1.5cm}
\newcommand\Overhang[1]{%
  \llap{\makebox[\TitleOverhang][l]{#1}}%
}
\titleformat{\section}
  {\normalfont\Large\bfseries}{\Overhang{\thesection}}{0em}{}
\titleformat{\subsection}
  {\normalfont\large\bfseries}{\Overhang{\thesubsection}}{0em}{}
\titleformat{\subsubsection}
  {\normalfont\normalsize\bfseries}{\Overhang{\thesubsubsection}}{0em}{}






\begin{document}

% % % % % % % % % CAPA % % % % % % % % % %
\begin{titlepage}
	\begin{center}
		Universidade Estadual de Campinas\\
		Faculdade de Engenharia Elétrica e de Computação\\ 
		Departamento de Engenharia de Computação e Automação Industrial \\ 
         \begin{figure}[htb]
             \centering
             \includegraphics[width=.08\linewidth]{imgs/unicamp.eps}
         \end{figure}
        \vspace{40pt}
        %\textbf{\Large{Sistema de realidade virtual para apresentação de conteúdo na língua de sinais brasileira}}\\
        \huge\textit{A voz na letra: \\ avaliação de representações visuais \\ de prosódia em tipografia}\\
        
	\end{center}
	
	\vspace{60pt}
	
    \hspace{0.5\textwidth}\begin{minipage}{0.5\textwidth}
        \begin{flushleft}     
    
            \textbf{\textsc{pesquisadores responsáveis}} \\ Profª Drª Paula Dornhofer Paro Costa (orientadora) \\  \vspace{0.2cm} Caluã de Lacerda Pataca (mestrando)
    
            \vspace{0.8cm}
    
            \textbf{\textsc{finalidade}} \\ Tese de mestrado
    
            \vspace{0.8cm}
    
            \textbf{\textsc{local}} \\ Universidade Estadual de Campinas \\  \vspace{0.2cm} Faculdade de Engenharia Elétrica e de~Computação \\  \vspace{0.2cm} Deptº de Engenharia de Computação e~Automação Industrial
               
               

    
    
        \end{flushleft}
    \end{minipage}	

    \vspace{2.4cm}     

    \begin{center}
        Apresentado em março de 2019
    \end{center}
	
% 	\vspace{80pt}
% 	\begin{flushleft}
%         \large{\textbf{Pesquisadores Responsáveis:}}\\
%         Profa. Dra. Paula Dornhofer Paro Costa (Orientadora)\\	
%         Caluã de Lacerda Pataca (Mestrando)\\
        
        
%         \vspace{10pt}
%         \large{\textbf{Finalidade:} \\Tese de Mestrado}\\
%         \vspace{10pt}
%         \large{\textbf{Local:} \\
%         Universidade Estadual de Campinas (Unicamp)
%         \\Faculdade de Engenharia Elétrica e de Computação (\textsc{feec})}\\
%         Depto. de Engenharia de Computação e Automação Industrial (\textsc{dca})\\
%         (e ambiente web?)
%         \vspace{15pt}        
% 	\end{flushleft}
% 	\vspace{3cm}
% 	\begin{center}
% 	 Data de Apresentação: Março, 2019
% 	\end{center}


\end{titlepage}

% % % % % % % % % % % % % % % % % % % % % % % % % % % % % % % % % %

% % % % % % % % % PRIMEIRA PÁGINA % % % % % % % % % % % % % % % % %
\begin{titlepage}
%\vspace{2cm}
% \begin{center}
\noindent \Large{\textbf{Dados dos Pesquisadores}}\\
% \end{center}

\vspace{10pt}
\begin{flushleft}
\large{\textbf{Nome}: Paula Dornhofer Paro Costa}\\
\large{\textbf{Função:} Pesquisadora Principal e Orientadora\\
\large{\textbf{Local de Trabalho:} Faculdade de Eng. Elétrica e de Computação (\textsc{feec}), \\ Depto. de Eng. de Computação e Automação Industrial (\textsc{dca})}\\
\large{\textbf{Endereço:} \\Av. Albert Einstein, 400, Cidade Universitária ``Zeferino Vaz'', Distrito de Barão Geraldo, Campinas – \textsc{sp}, \textsc{cep} 13083-852}\\
\large{\textbf{Telefone de Contato:} (19) 3521-0246 }\\
\large{\textbf{Currículo Lattes:} \url{http://lattes.cnpq.br/4518009815956207} }\\
\end{flushleft}

\vspace{10pt}

\begin{flushleft}
\large{\textbf{Nome}: Caluã de Lacerda Pataca }\\
\large{\textbf{Função:} Aluno de Mestrado}\\
\large{\textbf{Local de Trabalho:} Faculdade de Eng. Elétrica e de Computação (\textsc{feec}), \\ Depto. de Eng. de Computação e Automação Industrial (\textsc{dca})}\\
\large{\textbf{Endereço:} \\Av. Albert Einstein, 400, Cidade Universitária ``Zeferino Vaz'', Distrito de Barão Geraldo, Campinas – \textsc{sp}, \textsc{cep} 13083-852}\\
\large{\textbf{Telefone de Contato:} (19) 99289-1963 }\\
\large{\textbf{Currículo Lattes:} \url{http://lattes.cnpq.br/9013966339541270} }\\
\end{flushleft}
\vspace{10pt}

\vspace{10pt}

\end{titlepage}

\newpage

\abstract

\noindent A~leitura fluente se dá quando o leitor consegue relacionar internamente letras e símbolos com os sons da língua que ele em geral já traz consigo, em uma espécie de \textit{fala} interna. Em certos contextos esse processo funciona mal — algumas crianças em processo de alfabetização penam para se tornar leitoras fluentes, falantes de língua estrangeira nem sempre mapeiam bem som a letra, surdos não têm as referências sonoras etc. Nosso projeto se insere no rol de pesquisas em tecnologia tipográfica que visam manipular a forma gráfica do texto, visando complementar o conteúdo explícito da escrita com aspectos próprios da fala, investigando os impactos dessas intervenções sobre a leitura. Especificamente, construirá um algoritmo computacional que abstraia características acústicas da fala para mapeá-las visualmente em formas tipográficas, imbuindo no texto elementos expressivos da voz, em especial aqueles relacionados à prosódia. Nossa hipótese é que, na leitura de um texto composto com essa tipografia modulada pela voz, a relação entre intervenções gráficas e áudio será suficientemente intuitiva a ponto de um leitor qualquer saber apontar o vínculo. Adicionalmente, imaginamos que haverá com essas intervenções um aumento em medidas subjetivas relacionadas à \textit{imersão} do leitor em uma dada obra. A avaliação se dará em ambiente web e compreende duas etapas: Na primeira, os participantes serão expostos a gravações contendo falas de atores e, ao mesmo tempo, a textos com as transcrições dessas mesmas falas e em cujas formas foi feita a manipulação da tipografia por nosso algoritmo de extração e representação de prosódia. Espera-se então que o participante consiga relacionar o áudio a essas manipulações no texto. Na segunda etapa, haverá um questionário que buscará mensurar o impacto dessas manipulações no texto sobre a percepção subjetiva de seu conteúdo.

\newpage


\thispagestyle{empty}
\newpage


\newpage

\section*{Relevância Social}

\noindent A literatura documenta indícios empíricos de que as \emph{closed captions} em mídias audiovisuais trazem benefícios de compreensão, atenção e memória para diversos públicos, tais como: crianças, adultos e idosos; leitores experientes ou em fase de aprendizado; falantes ou não da língua em questão \cite{fiske_video_2015}. Além disso, a presença de \emph{closed captions} em programas de televisão é capaz de promover um entendimento mais aprofundado do conteúdo, em especial o conteúdo \emph{afetivo}, por crianças surdas e alfabetizadas\cite{murphy-berman_impact_1983}.

O presente projeto propõe o aperfeiçoamento nas tecnologias de representação de texto, por meio de técnicas que busquem representar qualidades expressivas presentes no registro sonoro e não necessariamente no registro visual. Os resultados esperados do trabalho têm o potencial de aumentar a acessibilidade e fruição de mídias audiovisuais por públicos diversos. Aplicações em potencial incluem o auxílio ao processo de alfabetização, o ensino de línguas estrangeiras e como tecnologia assistiva.

%Em cima dessa constatação emerge que aperfeiçoamentos nas tecnologias de representação de texto, em especial através de técnicas que busquem representar qualidades expressivas presentes no registro sonoro e não necessariamente no registro visual — sendo, assim, podem aumentar a acessibilidade e fruição de mídias audiovisuais.


%Aperfeiçoamentos nas tecnologias de representação de texto, em particular aquelas que buscam representar qualidades expressivas presentes no registro sonoro, têm o potencial de aumentar a acessibilidade e fruição de mídias audiovisuais. Esta hipótese


%De acordo com o Censo Demográfico de 2010 \cite{censo2010}, quase 9,7 milhões de cidadãos Brasileiros possuem alguma perda auditiva permanente. Há evidências empíricas \cite{murphy-berman_impact_1983} de que para crianças surdas e alfabetizadas, a presença de \textit{closed captions} em programas de televisão se traduz em um entendimento mais aprofundado do conteúdo, em especial em relação a seu conteúdo \textit{afetivo}.


%Mas os benefícios podem ser mais amplos: há também vasta literatura \cite{fiske_video_2015} documentando indícios empíricos de que as \textit{closed captions} trazem benefícios de compreensão, atenção e memória para diversos outros públicos: crianças, adultos e idosos; leitores experientes ou em fase de aprendizado; falantes ou não da língua em questão; etc. Assim, aperfeiçoamentos nessas legendas podem trazer benefícios a amplas parcelas da população.

%\todo{Paula, uma questão aqui: estou falando da possível importância do estudo para a população surda, mas *não* vou testar com surdos. Isso você acha que é um problema? Minha lógica, que pode ou não ter ficado clara no texto, é que as legendas são boas para qualquer público, o que significa que podemos testar com público em geral a princípio. }


\newpage


\thispagestyle{empty}
\newpage
\pagenumbering{arabic}

% % % % % % % % % % % % % % % % % % % % % % % % % % %

\renewcommand{\contentsname}{Sumário}
\tableofcontents

\newpage

%%%%%%%%%%
\section{Introdução}

\noindent Ler é uma habilidade cognitiva de alto nível, que exige longo período de treinamento e que envolve intenso processo neurológico. Dentre as estruturas cerebrais envolvidas no processo de leitura, é surpreendente notar que, além daquelas encarregadas do processamento de imagens, a leitura exige também o amplo emprego de estruturas que lidam com processamento de sons \cite[cap.7]{seidenberg2017}. Isso ocorre porque, mesmo quando lê silenciosamente, o leitor deduz (ou cria) em sua voz interna a musicalidade do texto, habilidade fundamentalmente relacionada à compreensão e interpretação do mesmo.

Investigar como essa ``voz'' emerge a partir do texto não se trata de uma questão meramente \textit{estética}. Certos tipos de dislexia, por exemplo, parecem antes causados por problemas nessas estruturas neurológicas que processam sons do que deficiências no processamento de imagens, mesmo que se manifestem sob a forma de dificuldades na  leitura \cite[cap.8]{seidenberg2017}. Além disso, crianças em processo de alfabetização que leem de maneira monótona, ou seja, que não conseguem extrair do texto a expressividade natural da fala, tendem a desenvolver problemas de compreensão \cite{bessemans2017}. Finalmente — e ao contrário da noção vendida por certos cursos de leitura dinâmica de que uma leitura sem subvocalização traria ganhos de velocidade sem perdas na compreensão —, o leitor experiente se vale dessa voz interna para ter acesso à informação fonológica contida no texto e, assim, reduzir ambiguidades e facilitar a compreensão \cite[cap.4]{seidenberg2017}.

Neste contexto, o presente projeto se insere no âmbito da pesquisa em tecnologia tipográfica. Se debruça sobre algoritmos computacionais de transformação visual de texto almejando transformar o processo de aquisição de informação por meio da leitura mais intuitivo e acessível, com potenciais aplicações no auxílio ao processo de alfabetização, ensino de línguas estrangeiras e, entre ainda outras, como tecnologia assistiva.

Em particular, a presente pesquisa propõe investigar e avaliar algoritmos de mapeamento automático de parâmetros acústicos da fala expressiva, especialmente aqueles relacionados à expressão da prosódia, para parâmetros tipográficos do texto.



%%%%%%%%%%
\section{Objetivos}

\noindent Este projeto tem como objetivo avaliar um modelo de mapeamento de características acústicas da voz para modulações visuais em uma família tipográfica. 
São objetivos específicos do projeto:
\begin{itemize}
    \item medir quão bem um participante qualquer consegue inferir a ligação entre fala e texto modificado;
    \item apurar que efeitos essa abordagem pode ter na percepção subjetiva de um participante em relação a um dado conteúdo.
\end{itemize} 

%%%%%%%%%%
\section{Hipóteses}

\noindent A primeira hipótese a ser testada é que nossa abordagem para mapear a prosódia contida na voz em modificações visuais tipográficas produzirá resultados suficientemente auto-explicativos a ponto de que um participante da avaliação, não tendo recebido treinamento prévio, consiga identificar que um dado texto corresponde a uma dada enunciação sonora do mesmo.

A segunda hipótese a ser testada é que essas modificações causarão um aumento na imersão do participante em uma dada obra audiovisual, medida através do aumento em medidas subjetivas auto-reportadas como \textit{transporte}, \textit{identificação com as personagens} e \textit{realismo percebido}.

%%%%%%%%%%
\section{Local de Realização da Pesquisa}

\noindent Os participantes serão divididos em dois grupos caracterizados pelo ambiente de realização de pesquisa, como se segue:

\begin{itemize}
    \item \textbf{Grupo Avaliado em Condições Controladas}: Os participantes deste grupo participarão de avaliação perceptual nas dependências da Faculdade de Engenharia Elétrica e de Computação (FEEC) da Universidade Estadual de Campinas (Unicamp) ou, alternativamente, em salas de aula ou laboratórios da Unicamp previamente agendados para a avaliação. As características principais deste grupo estão associadas ao controle das condições do ambiente de avaliação que garantirá a privacidade dos participantes durante a avaliação, deverá ter iluminação adequada, cadeiras confortáveis e ser suficientemente isolado de ruídos externos que distraiam ou atrapalhem a audição dos estímulos em áudio que serão apresentados aos participantes. Além disso, neste grupo as avaliações serão sempre realizadas sempre sob a supervisão dos pesquisadores, mediante agendamento. 
    \item \textbf{Grupo Avaliado Online}: Este grupo é formado pelo conjunto de participantes que optarem por realizar a avaliação online em ambientes livremente escolhidos pelos participantes e sem controle das condições ambiente de conforto, ruído e iluminação. Este segundo grupo terá recebido um convite por e-mail contendo uma apresentação do estudo, a indicação do período em que o teste estará disponível online, o perfil esperado para participação e os requisitos mínimos computacionais necessários para acessar o site. Ficam assim livres para escolher local e horário que melhor lhes convir para realização da avaliação.
\end{itemize}

%As avaliações serão realizadas em dois contextos: em um, sob supervisão dos pesquisadores nas dependências do Laboratório do Departamento de Engenharia de Computação e Automação Industrial (\textsc{lca}), na Faculdade de Engenharia Elétrica e de Computação (\textsc{feec}); em outro, em ambientes livremente escolhidos pelos participantes que optem por fazer o experimento em sua modalidade online. 

%Este segundo grupo terá recebido um convite por e-mail contendo uma apresentação do estudo, a indicação do período em que o teste estará disponível online, o perfil esperado para participação e os requisitos mínimos computacionais necessários para acessar o site. Ficam assim livres para escolher local e horário que melhor lhes convir para realização do experimento.

%As avaliações presenciais do primeiro grupo, de controle para o segundo, serão realizadas nas dependências do \textsc{lca}, na \textsc{feec}, em horários previamente agendados de comum acordo entre os pesquisadores responsáveis e o participante voluntário.

%Os participantes serão informados durante o recrutamento e explicitamente via Termo de Consentimento Livre e Esclarecido (\textsc{tcle}) que o projeto não arcará com eventuais custos de deslocamento até o local da avaliação. Os agendamentos serão realizados de modo que coincidam com dia e horários de atendimento usual dos participantes pelo \textsc{lca}.

A sala onde será realizada os testes pode ser caracterizada como uma sala contendo computador, monitor e cadeiras e, no momento da realização dos testes estará dedicada à avaliação a fim de evitar distrações para os participantes e garantir a privacidade dos mesmos durante a avaliação.

%%%%%%%%%%
\section{População a ser estudada}

\noindent Serão realizadas chamadas de recrutamento para indivíduos fluentes em português, alfabetizados, maiores de idade e, para aqueles que decidirem fazer a avaliação em sua modalidade online, com acesso a um computador com sistema de som, internet e versões recentes dos navegadores mais comuns.

%%%%%%%%%%
\section{Garantias éticas aos participantes da pesquisa}
\noindent A metodologia empregada pelo estudo e seus procedimentos estarão descritos no Termo de Consentimento Livre e Esclarecido (\textsc{tcle}), em linguagem clara e acessível, o qual será apresentado a cada participante. Para os participantes que optarem por fazer o teste na modalidade online, a avaliação só se iniciará a partir do momento em que sinalizarem que leram e que concordam com o \textsc{tcle}, que estará disponível para leitura e \textit{download} no site da avaliação. 

Para os participantes do grupo de controle, que farão a avaliação sob supervisão em ambiente controlado, a anuência aos termos apresentados será registrada por meio da assinatura o termo em duas vias, sendo que uma via ficará em mãos dos pesquisadores e a outra com o participante.

Os dados serão salvos de maneira anônima: todas informações coletadas para delimitar o perfil do participante (idade, sexo, formação e área de atuação) estarão associadas a suas respostas, mas não haverá nenhum vínculo entre esses dados e informações que identifiquem individualmente cada participante.

%%%%%%%%%%


%%%%%%%%%%
\section{Metodologia Proposta}

\noindent Como já exposto, a avaliação terá duas configurações: um primeiro grupo realizará o teste em ambiente controloado e sob supervisão dos pesquisadores; o segundo, em ambiente livre de acordo com a conveniência do participante, demandando apenas um computador com requisitos pré-determinados.

\hl{Para os participantes que farão a avaliação sob supervisão, antes do início do teste propriamente o \textsc{tcle} será apresentado verbalmente e em duas vias por escrito, as quais o participante deverá ler e assinar para prosseguir, ficando com uma das vias.}

\todo{Paula, estou imaginando que para o grupo de controle isso aqui teria que ser físico, mas se você achar que rola ser exclusivamente via site: melhor!}

\textcolor{blue}{Eu acho melhor sermos conservadores, vamos coletar assinaturas físicas.}

%As sessões supervisionadas ocorrerão em uma sala no Laboratório do Departamento de Engenharia de Computação e Automação Industrial (\textsc{lca}).

Nem os participantes do grupo sob supervisão tampouco os demais serão filmados ou fotografados durante a avaliação.

Em ambos os casos, a avaliação consiste no participante interagir com uma página web, seguindo os seguintes passos:

\begin{itemize}
    \item Apresentação resumida do projeto por escrito;
    \item Para os participantes que farão a avaliação na modalide online, apresentação do Termo de Consentimento Livre e Esclarecido (\textsc{tcle}). Para prosseguir com a avaliação o participante deverá declarar ter lido e que está de acordo com os termos propostos;
    \item Questionário inicial para coleta de dados demográficos — idade e grau de formação — resposta obrigatória para prosseguir —, e sexo e área de trabalho e/ou estudo — resposta opcional. Para os participantes que farão a avaliação na modalidade online, nesta etapa serão também coletados automaticamente dados sobre localização de acesso e informações sobre o computador e navegador utilizados (sistema operacional, tamanho de tela, versões etc). Não serão coletadas informações que permitam recuperar a identidade do participante, cujas respostas e dados serão salvos sempre de maneira anônima; 
    \item Início do teste com áudio e tipografia. Em cada rodada, será apresentado um áudio contendo uma fala e sua transcrição. Em alguns casos, essa transcrição terá tido o desenho de suas letras modificado de acordo com nossa abordagem de extração e representação de prosódia. Em outros casos, essas modificações não corresponderão ao áudio. O participante deverá relacionar a tipografia que corresponde ao áudio, produzindo um índice de acerto que nos indicará a efetividade de nossa abordagem quando comparada com uma tipografia com modulações não correspondentes ao som ouvido;
    \item Na segunda etapa haverá um conjunto de perguntas do tipo ``escala de Likert'' com sete gradações, adaptadas de \citeonline{translatology}, visando mensurar medidas subjetivas como \textit{transporte} (e.g. “enquanto fazia o teste da etapa anterior, quão bem você manteve consciência do ambiente em seu entorno?”), \textit{identificação com as personagens} (e.g. “quão claras estavam as percepções e emoções das personagens para você?”) e \textit{realismo percebido} (e.g. “quão realistas lhe pareceram as personagens?”). \todo{Paula, note que esta etapa não pensei para o primeiro teste e sim para o seguinte, com legendas. Mas eis o que me ocorreu: talvez possamos jogar umas questões de escala Likert mesmo no primeiro teste e correlacioná-las aos índices de performance dos participantes. Não responde ao nosso objetivo principal, mas não o influencia negativamente (seria ao final) e talvez traga algum elemento novo.}
\end{itemize}

Tanto no teste a ser realizado na \textsc{feec} quanto na modalidade online, o participante pode interromper a avaliação a qualquer momento, sem incorrer em qualquer prejuízo a si.

\todo{Paula: não estou falando sobre o conteúdo das frases para que possa usar um vídeo no segundo experimento (o das legendas) sem ter que escolher um filme agora. Você acha isso um problema?}
\textcolor{blue}{Perfeito. Como eu disse, vamos colocar o maior número de ideias possíveis nesse projeto. Fazer menos do que foi dito que faria não seria um problema (ao contrário de fazer mais do que foi inicialmente contemplado).}

% As animações são frases sinalizadas em Libras traduzidas do livro de ciências do Ensino Fundamental I. As frases das duas primeiras etapas serão diferentes, além disso, para cada participante, o conjunto de frases será escolhido aleatoriamente.


% Likert
%As opções de respostas, concordância ou não para cada frase, serão:
%\begin{enumerate}
%\item Discordo totalmente
%\item Discordo parcialmente
%\item Indiferente
%\item Concordo parcialmente
%\item Concordo totalmente
%\end{enumerate}


\subsection{Plano de Recrutamento}

\noindent Recrutaremos participantes que cumpram os critérios de inclusão através de mensagem de e-mail e divulgação em redes sociais direcionadas a alunos de graduação, pós-graduação e funcionários da Unicamp.

Adicionalmente, serão afixados cartazes em quadros de aviso pela Unicamp.

%%%%%%%%%%
\section{Cronograma}
\noindent O projeto tem previsão de duração de aproximadamente 4 meses, com recrutamento com início previsto após a aprovação do projeto no Comitê de Ética em Pesquisa, seguindo um plano de trabalho divido nas etapas apresentadas na Tabela~\ref{table1}, conforme cronograma apresentado.


\begin{table}[H]
\centering
\resizebox{\textwidth}{!}{%
\begin{tabular}{lll}
\textbf{Descrição Atividade} & \textbf{Período Planejado} & \textbf{Observações}  \\ \cline{1-3}
Desenvolvimento do site & 11/3/19 — 21/4/19 & \begin{tabular}[l]{@{}l@{}}Esta etapa refere-se ao desenvolvimento \\ do software e não envolve a coleta de dados.\end{tabular} \\ \hdashline[0.5pt/3pt]
Recrutamento & 22/4/19 — 28/4/19 & \begin{tabular}[l]{@{}l@{}}Etapa de recrutamento dos \\ participantes da pesquisa.\end{tabular} \\ \hdashline[0.5pt/3pt]
Execução do Estudo & 29/4/19 — 26/5/19 & \begin{tabular}[l]{@{}l@{}}Avaliações supervisionadas e \\ período em que o site estará aberto.\end{tabular} \\ \hdashline[0.5pt/3pt]
Análise dos Resultados & 27/5/19 — 17/6/19 & \begin{tabular}[l]{@{}l@{}}Etapa de compilação \\ e análise dos resultados.\end{tabular} \\ \hdashline[0.5pt/3pt]
Publicação do Artigo & 18/6/19 — 15/7/19 & \begin{tabular}[l]{@{}l@{}}Etapa de registro dos resultados obtidos e\\ e escrita de artigo para divulgação dos \\ principais resultados.\end{tabular} \\ \hdashline[0.5pt/3pt]
\end{tabular}%
}
\caption{Cronograma do Projeto}
\label{table1}
\end{table}

\noindent O recrutamento para a participação na pesquisa só será iniciada efetivamente após aprovação do projeto pelo Comitê de Ética.

%%%%%%%%%%
\section{Orçamento}
\noindent A hospedagem do site será feita em servidores da empresa de infraestrutura web \textit{Surge}, que apresenta um plano gratuito sob o qual nossa demanda prevista se enquadra. O banco de dados com as respostas dos participantes ficará no serviço \textit{Firebase}, da empresa \textit{Google}, que também possui planos gratuitos para sites de menor complexidade, como o nosso.

Os recursos necessários para a realização da parte presencial e supervisionada da pesquisa, tais como mesas e computador, serão emprestados da Faculdade de Engenharia Elétrica e de Computação. Porém, teremos as despesas suplementares referente a pesquisa que estão detalhadas abaixo:

\begin{table}[h]
\caption{Detalhamento do Orçamento}
\centering
\vspace{0.5cm}
\begin{tabular}{l|r}
Identificação do Orçamento & Valor em Reais (R\$) \\
\hline 
Água        & 10,00 \\
Impressão de Instruções e Termos  & 15,00 \\
Impressão de Cartazes  & 25,00 \\
\hline 
Total & 50,00
\end{tabular}
\end{table}

%%%%%%%%%%
\section{Critérios de inclusão e exclusão dos participantes da pesquisa}

\noindent Para ser considerado elígivel o participante deverá ser fluente em português, alfabetizado e maior de idade. Adicionalmente, os participantes da modalidade online deverão ter acesso a um computador com sistema de som, internet e versões recentes dos navegadores mais comuns.

Serão excluídos deste estudo indivíduos:

\begin{itemize}
\item indivíduos em situação de vulnerabilidade como pacientes em tratamento de doenças graves ou com problemas mentais;
\item esteja com algum tipo de debilidade no seu senso de movimento e equilíbrio (dores de cabeça, gripe, enxaquecas, ansiedade ou estar sob a influência de álcool e drogas);
\item possua alguma deficiência visual que impeça o reconhecimento de imagens com detalhes finos.
\end{itemize}

Distúrbios ou doenças oculares como daltonismo, ametropias ou nistagmo não são critérios excludentes, mas podem dificultar a interpretação dos estímulos visuais durante a pesquisa. Serão aceitos indivíduos que utilizam ou não lentes corretivas. 

O local onde a parte supervisionada dos estudos será realizada possui acesso para cadeirantes.

%%%%%%%%%%
\section{Riscos e Benefícios envolvidos na execução da pesquisa}
\noindent Tanto os participantes que fizerem a avaliação na modalidade online quanto aqueles que o fizerem no \textsc{lca} da \textsc{feec} serão expostos aos áudios e imagens e responderão aos questionários em uma única sessão.

Todos os equipamentos utilizados nesta avaliação são seguros e indolores. 

Não há riscos previsíveis para os participantes da pesquisa.

Para garantir conforto e evitar cansaço, os participantes do grupo de controle terão acesso a cadeiras confortáveis e o ambiente estará adequadamente iluminado e climatizado. Além disso, os participantes serão monitorados frequentemente e, a qualquer momento, o usuário poderá sinalizar qualquer desconforto e a avaliação será interrompida imediatamente. 

Caso possíveis sintomas apresentados persistam após a interrupção da avaliação, o participante será encaminhado ao atendimento de saúde disponível no Pronto Socorro do Hospital das Clínicas da Unicamp.%a 5 minutos do \textsc{lca}.}

Não existem benefícios previstos para os participantes da pesquisa de forma individual. Porém, o resultado desta pesquisa visa trazer benefícios a todos que consomem mídias audiovisuais com \textit{closed captions} ou legendas.

%%%%%%%%%%
\section{Critérios de Encerramento ou Suspensão da Pesquisa}
\noindent O encerramento ou suspensão de exibição dos áudios e textos se dará caso o participante sofra qualquer tipo de desconforto ou decida encerrar a avaliação, sem que seja necessário nenhum motivo específico. 

%%%%%%%%%%
\section{Resultados do Estudo}
\noindent O estudo tem como resultado final a apresentação dos resultados de compreensão da associação entre áudio e modificações visuais na tipografia correspondentes, bem como de o quanto essa abordagem modifica medidas subjetivas auto-reportadas de engajamento com conteúdos audiovisuais.

%%%%%%%%%%
\section{Divulgação dos Resultados}
\noindent A divulgação deste estudo se dará em forma de artigo científico. Além disso, os resultados deste estudo farão parte de dissertação de mestrado na Faculdade de Engenharia Elétrica e de Computação (\textsc{feec}).

%%%%%%%%%%
\section{Requisitos Específicos do Protocolo de Pesquisa}
\noindent Não foram identificados outros requisitos específicos do protocolo de pesquisa além dos já mencionados nas seções anteriores.

\addcontentsline{toc}{section}{Referências Bibliográficas}

% \bibliographystyle{plain}  %% nome-ano %abnt
\bibliographystyle{abntex2-alf}
\bibliography{ref}



\newpage
\addcontentsline{toc}{section}{Anexo}
%\section*{Anexo}

%\subsection{Questões}

\end{document}


